\begin{thebibliography}{99}
\footnotesize\itemsep=0pt

\bibitem{BCDX2018}
Bai L., Chen X., Ding M., Xu F., A quantum analog of generalized cluster
  algebras, \href{https://doi.org/10.1007/s10468-017-9743-7}{\textit{Algebr. Represent. Theory}} \textbf{21} (2018), 1203--1217,
  \href{https://arxiv.org/abs/1610.09803}{arXiv:1610.09803}.

\bibitem{BCDX2023}
Bai L., Chen X., Ding M., Xu F., Generalized quantum cluster algebras: the
  {L}aurent phenomenon and upper bounds, \href{https://doi.org/10.1016/j.jalgebra.2022.12.012}{\textit{J.~Algebra}} \textbf{619}
  (2023), 298--322, \href{https://arxiv.org/abs/2203.06928}{arXiv:2203.06928}.

\bibitem{Berenstein_Zelevinsky2005}
Berenstein A., Zelevinsky A., Quantum cluster algebras, \href{https://doi.org/10.1016/j.aim.2004.08.003}{\textit{Adv. Math.}}
  \textbf{195} (2005), 405--455, \href{https://arxiv.org/abs/math.QA/0404446}{arXiv:math.QA/0404446}.

\bibitem{Chekhov_Shapiro}
Chekhov L., Shapiro M., Teichm\"uller spaces of {R}iemann surfaces with
  orbifold points of arbitrary order and cluster variables,  \href{https://doi.org/10.1093/imrn/rnt016}{\textit{Int. Math.
  Res. Not.}} \textbf{2014} (2014), 2746--2772, \href{https://arxiv.org/abs/1111.3963}{arXiv:1111.3963}.

\bibitem{Davison18}
Davison B., Positivity for quantum cluster algebras,  \href{https://doi.org/10.4007/annals.2018.187.1.3}{\textit{Ann. of Math.}}
  \textbf{187} (2018), 157--219, \href{https://arxiv.org/abs/1601.07918}{arXiv:1601.07918}.

\bibitem{fock2009}
Fock V.V., Goncharov A.B., Cluster ensembles, quantization and the dilogarithm,
   \href{https://doi.org/10.24033/asens.2112}{\textit{Ann. Sci. \'Ec. Norm. Sup\'er.}} \textbf{42} (2009), 865--930,
  \href{https://arxiv.org/abs/math.AG/0311245}{arXiv:math.AG/0311245}.

\bibitem{fomin_zelevinsky_2002}
Fomin S., Zelevinsky A., Cluster algebras.~{I}. {F}oundations, \href{https://doi.org/10.1090/S0894-0347-01-00385-X}{\textit{J.~Amer.
  Math. Soc.}} \textbf{15} (2002), 497--529, \href{https://arxiv.org/abs/math.RT/0104151}{arXiv:math.RT/0104151}.

\bibitem{fomin_zelevinsky_2007}
Fomin S., Zelevinsky A., Cluster algebras.~{IV}. {C}oefficients,
  \href{https://doi.org/10.1112/S0010437X06002521}{\textit{Compos. Math.}} \textbf{143} (2007), 112--164,
  \href{https://arxiv.org/abs/math.RA/0602259}{arXiv:math.RA/0602259}.

\bibitem{Gleitz2015}
Gleitz A.-S., Generalised cluster algebras and {$q$}-characters at roots of
  unity, in Proceedings of {FPSAC} 2015, \textit{Discrete Math. Theor. Comput. Sci. Proc.}, \href{https://doi.org/10.46298/dmtcs.2479}{Discrete Mathematics \& Theoretical Computer Science (DMTCS)}, Nancy,
  2015, 357--368.

\bibitem{GHKK2018}
Gross M., Hacking P., Keel S., Kontsevich M., Canonical bases for cluster
  algebras,  \href{https://doi.org/10.1090/jams/890}{\textit{J.~Amer. Math. Soc.}} \textbf{31} (2018), 497--608,
  \href{https://arxiv.org/abs/1411.1394}{arXiv:1411.1394}.

\bibitem{gupta2019}
Gupta M., A formula for {$F$}-polynomials in terms of $c$-vectors and
  stabilization of {$F$}-polynomials, \href{https://arxiv.org/abs/1812.01910}{arXiv:1812.01910}.

\bibitem{Iwaki-Nakanishi2016}
Iwaki K., Nakanishi T., Exact {WKB} analysis and cluster algebras {II}:
  {S}imple poles, orbifold points, and generalized cluster algebras,
  \href{https://doi.org/10.1093/imrn/rnv270}{\textit{Int. Math. Res. Not.}} \textbf{2016} (2016), 4375--4417,
  \href{https://arxiv.org/abs/1409.4641}{arXiv:1409.4641}.

\bibitem{Keller2012}
Keller B., Cluster algebras and derived categories, in Derived {C}ategories in
  {A}lgebraic {G}eometry, \textit{EMS Ser. Congr. Rep.},  \href{https://doi.org/10.4171/115-1/6}{European Mathematical Society (EMS)}, Z\"urich, 2013, 123--183, \href{https://arxiv.org/abs/1202.4161}{arXiv:1202.4161}.

\bibitem{KD2020}
Keller B., Demonet L., A survey on maximal green sequences, in Representation
  {T}heory and {B}eyond, \textit{Contemp. Math.}, Vol.~758, \href{https://doi.org/10.1090/conm/758/15239}{American
  Mathematical Society}, Providence, RI, 2020, 267--286, \href{https://arxiv.org/abs/1904.09247}{arXiv:1904.09247}.

\bibitem{LFM18}
Labardini-Fragoso D., Mou L., Gentle algebras arising from surfaces with
  orbifold points of order~3, Part~{I}: scattering diagrams, \href{https://doi.org/10.1007/s10468-023-10233-x}{\textit{Algebr.
  Represent. Theory}} \textbf{27} (2024), 679--722, \href{https://arxiv.org/abs/2203.11563}{arXiv:2203.11563}.

\bibitem{LF-Velasco2018}
Labardini-Fragoso D., Velasco D., On a family of {C}aldero--{C}hapoton algebras
  that have the {L}aurent phenomenon, \href{https://doi.org/10.1016/j.jalgebra.2018.11.012}{\textit{J.~Algebra}} \textbf{520} (2019),
  90--135, \href{https://arxiv.org/abs/1704.07921}{arXiv:1704.07921}.

\bibitem{LMN2023}
Lin F., Musiker G., Nakanishi T., Two formulas for {$F$}-polynomials,
   \href{https://doi.org/10.1093/imrn/rnad074}{\textit{Int. Math. Res. Not.}} \textbf{2024} (2024), 613--634,
  \href{https://arxiv.org/abs/2112.11839}{arXiv:2112.11839}.

\bibitem{Nakanishi-qgca}
Nakanishi T., Quantum generalized cluster algebras and quantum dilogarithms of
  higher degrees, \href{https://doi.org/10.1007/s11232-015-0377-9}{\textit{Theoret. and Math. Phys.}} \textbf{185} (2015),
  1759--1768, \href{https://arxiv.org/abs/1410.0584}{arXiv:1410.0584}.

\bibitem{Nakanishi14}
Nakanishi T., Structure of seeds in generalized cluster algebras,
  \href{https://doi.org/10.2140/pjm.2015.277.201}{\textit{Pacific~J. Math.}} \textbf{277} (2015), 201--217, \href{https://arxiv.org/abs/1409.5967}{arXiv:1409.5967}.

\bibitem{nakanishi2023}
Nakanishi T., Cluster algebras and scattering diagrams, \textit{MSJ Memoirs}, Vol.~41, \href{https://doi.org/10.1142/e073}{Mathematical Society of Japan}, Tokyo, 2023, \href{https://arxiv.org/abs/2201.11371}{arXiv:2201.11371}.

\bibitem{Thao2011}
Tran T., {$F$}-polynomials in quantum cluster algebras,  \href{https://doi.org/10.1007/s10468-010-9226-6}{\textit{Algebr.
  Represent. Theory}} \textbf{14} (2011), 1025--1061, \href{https://arxiv.org/abs/0904.3291}{arXiv:0904.3291}.

\end{thebibliography}\LastPageEnding
